% --- Archon physics formalization source begin ---
% archon:physics
% archon:covers PhyXMiniProblems/problem_phyx_mini_0209.lean
% archon:source-report reports/phyx_mini/problem_phyx_mini_0209.source.json

\chapter{Physics problem phyx\_mini\_0209}
\label{ch:PhyXMiniProblems_problem_phyx_mini_0209}

\paragraph{Problem source.}
\#\# Physical scenario

For any type of wave that reaches a boundary beyond which its speed is increased, there is a maximum incident angle if there is to be a transmitted refracted wave. This maximum incident angle \textbackslash{}( \textbackslash{}theta\_\{i\textbackslash{}mathrm\{M\}\} \textbackslash{}) corresponds to an angle of refraction equal to \textbackslash{}( 90\textasciicircum{}\{\textbackslash{}circ\} \textbackslash{}). If \textbackslash{}( \textbackslash{}theta\_i > \textbackslash{}theta\_\{i\textbackslash{}mathrm\{M\}\} \textbackslash{}), all the wave is reflected at the boundary and none is refracted, because refraction would correspond to \textbackslash{}( \textbackslash{}sin \textbackslash{}theta\_r > 1 \textbackslash{}) (where \textbackslash{}( \textbackslash{}theta\_r \textbackslash{}) is the angle of refraction), which is impossible. This phenomenon is referred to as total internal reflection.His voice is \textbackslash{}( 1.8\textbackslash{},\textbackslash{}mathrm\{m\} \textbackslash{} above the ground.The speed of sound is about \textbackslash{}( 343\textbackslash{},\textbackslash{}mathrm\{m/s\} \textbackslash{}) in air and \textbackslash{}( 1440\textbackslash{},\textbackslash{}mathrm\{m/s\} \textbackslash{}) in water.

\#\# Question

How far from the bank should a trout fisherman stand (figure) so trout won't be frightened by his voice?

\#\# Answer choices

- A: A: 0.56m

- B: B: 0.52m

- C: C: 0.48m

- D: D: 0.44m

\#\# Figure caption (auxiliary; use the image as primary evidence)

The image depicts an individual standing on the grassy bank of a body of water, using a fishing rod. The rod casts a line into the water, where a fish can be seen swimming.

Below the rod, a pink line is drawn at an angle, labeled "Sound wave." The wavefronts are shown as blue curves along this line. 

A dashed line is perpendicular to the water's surface, forming an angle with the sound wave, called "θ\_i." 

The horizontal distance from the person to the point where the sound wave hits the water is labeled as "x."

\#\# Dataset metadata

Wave Properties; Physical Model Grounding Reasoning, Spatial Relation Reasoning

\paragraph{Recorded answer/context.}
D: D: 0.44m

\paragraph{Figure/image path.}
/root/proposal\_for\_physic/science-mango/phyx\_mini\_run/phyx\_data/test\_image/209.png

\paragraph{Formalization target.}
create a compiling Lean file with sorry bodies at `PhyXMiniProblems/problem\_phyx\_mini\_0209.lean`. The Lean declarations must preserve the physical quantities, dimensions or dimensional roles, figure labels, governing-law hypotheses, and final relation expressed by this problem.
Use Mathlib/Physlib names found through LeanExplore where available. If a physics API is missing, introduce faithful local abstractions rather than scalar placeholder aliases.

\begin{theorem}[Physics formalization target]
\label{thm:physics:phyx_mini_0209:target}
\lean{PhyXMiniProblems.ProblemPhyXMini0209.critical_bank_setback_is_recorded_choice_D}
\uses{def:physics:phyx-mini-0209:phyxminiproblems-problemphyxmini0209-troutfishersoundsetup, def:physics:phyx-mini-0209:phyxminiproblems-problemphyxmini0209-matchesproblemdata, def:physics:phyx-mini-0209:phyxminiproblems-problemphyxmini0209-matchesprimaryfigure, def:physics:phyx-mini-0209:phyxminiproblems-problemphyxmini0209-hasphysicalparameters, def:physics:phyx-mini-0209:phyxminiproblems-problemphyxmini0209-satisfiescriticalacousticrefractionphysics, def:physics:phyx-mini-0209:phyxminiproblems-problemphyxmini0209-isminimumsafesetback, def:physics:phyx-mini-0209:phyxminiproblems-problemphyxmini0209-recordedanswerchoice, def:physics:phyx-mini-0209:phyxminiproblems-problemphyxmini0209-matchesdisplayedcentimeter}
The assigned autoformalize agent should translate this physics problem into Lean declarations in the covered file, with theorem and lemma proof bodies written as `by sorry`.
\end{theorem}
\begin{proof}
This is an autoformalization task, not a proof task. Produce statements that can later be proved by the physics prover without weakening the physical model.
\end{proof}

% --- Archon declaration topology begin ---
\section{Declaration topology}
\begin{definition}[Length Magnitude]
  \label{def:physics:phyx-mini-0209:phyxminiproblems-problemphyxmini0209-lengthmagnitude}
  \lean{PhyXMiniProblems.ProblemPhyXMini0209.LengthMagnitude}
  A nonnegative physical length, independent of the unit used to read it
\end{definition}
\begin{proof}
  This is the declared model definition of Length Magnitude.
\end{proof}
\begin{definition}[length In Meters]
  \label{def:physics:phyx-mini-0209:phyxminiproblems-problemphyxmini0209-lengthinmeters}
  \lean{PhyXMiniProblems.ProblemPhyXMini0209.lengthInMeters}
  \uses{def:physics:phyx-mini-0209:phyxminiproblems-problemphyxmini0209-lengthmagnitude}
  Read a physical length in metres
\end{definition}
\begin{proof}
  This is the declared model definition of length In Meters.
\end{proof}
\begin{definition}[speed In Meters Per Second]
  \label{def:physics:phyx-mini-0209:phyxminiproblems-problemphyxmini0209-speedinmeterspersecond}
  \lean{PhyXMiniProblems.ProblemPhyXMini0209.speedInMetersPerSecond}
  Read Physlib's nonnegative physical speed in metres per second
\end{definition}
\begin{proof}
  This is the declared model definition of speed In Meters Per Second.
\end{proof}
\begin{definition}[Acoustic Medium]
  \label{def:physics:phyx-mini-0209:phyxminiproblems-problemphyxmini0209-acousticmedium}
  \lean{PhyXMiniProblems.ProblemPhyXMini0209.AcousticMedium}
  The two propagation media separated by the horizontal water surface
\end{definition}
\begin{proof}
  This is the declared model definition of Acoustic Medium.
\end{proof}
\begin{definition}[Figure Object]
  \label{def:physics:phyx-mini-0209:phyxminiproblems-problemphyxmini0209-figureobject}
  \lean{PhyXMiniProblems.ProblemPhyXMini0209.FigureObject}
  Physical objects visible in the fisherman-and-trout figure
\end{definition}
\begin{proof}
  This is the declared model definition of Figure Object.
\end{proof}
\begin{definition}[Figure Object Role]
  \label{def:physics:phyx-mini-0209:phyxminiproblems-problemphyxmini0209-figureobjectrole}
  \lean{PhyXMiniProblems.ProblemPhyXMini0209.FigureObjectRole}
  Physical role assigned to each depicted object
\end{definition}
\begin{proof}
  This is the declared model definition of Figure Object Role.
\end{proof}
\begin{definition}[Figure Line]
  \label{def:physics:phyx-mini-0209:phyxminiproblems-problemphyxmini0209-figureline}
  \lean{PhyXMiniProblems.ProblemPhyXMini0209.FigureLine}
  The four distinguished lines or segments in the supplied image
\end{definition}
\begin{proof}
  This is the declared model definition of Figure Line.
\end{proof}
\begin{definition}[Figure Line Role]
  \label{def:physics:phyx-mini-0209:phyxminiproblems-problemphyxmini0209-figurelinerole}
  \lean{PhyXMiniProblems.ProblemPhyXMini0209.FigureLineRole}
  Geometric or acoustic role of a distinguished figure line
\end{definition}
\begin{proof}
  This is the declared model definition of Figure Line Role.
\end{proof}
\begin{definition}[Figure Angle Label]
  \label{def:physics:phyx-mini-0209:phyxminiproblems-problemphyxmini0209-figureanglelabel}
  \lean{PhyXMiniProblems.ProblemPhyXMini0209.FigureAngleLabel}
  The incidence-angle symbol printed beside the dashed normal
\end{definition}
\begin{proof}
  This is the declared model definition of Figure Angle Label.
\end{proof}
\begin{definition}[Figure Distance Label]
  \label{def:physics:phyx-mini-0209:phyxminiproblems-problemphyxmini0209-figuredistancelabel}
  \lean{PhyXMiniProblems.ProblemPhyXMini0209.FigureDistanceLabel}
  The horizontal-distance symbol printed below the figure
\end{definition}
\begin{proof}
  This is the declared model definition of Figure Distance Label.
\end{proof}
\begin{definition}[Trout Fisher Figure]
  \label{def:physics:phyx-mini-0209:phyxminiproblems-problemphyxmini0209-troutfisherfigure}
  \lean{PhyXMiniProblems.ProblemPhyXMini0209.TroutFisherFigure}
  \uses{def:physics:phyx-mini-0209:phyxminiproblems-problemphyxmini0209-figureobject, def:physics:phyx-mini-0209:phyxminiproblems-problemphyxmini0209-figureobjectrole, def:physics:phyx-mini-0209:phyxminiproblems-problemphyxmini0209-figureline, def:physics:phyx-mini-0209:phyxminiproblems-problemphyxmini0209-figurelinerole, def:physics:phyx-mini-0209:phyxminiproblems-problemphyxmini0209-figureanglelabel, def:physics:phyx-mini-0209:phyxminiproblems-problemphyxmini0209-figuredistancelabel}
  Qualitative labels and roles read directly from the primary raster
\end{definition}
\begin{proof}
  This is the declared model definition of Trout Fisher Figure.
\end{proof}
\begin{definition}[Trout Fisher Sound Setup]
  \label{def:physics:phyx-mini-0209:phyxminiproblems-problemphyxmini0209-troutfishersoundsetup}
  \lean{PhyXMiniProblems.ProblemPhyXMini0209.TroutFisherSoundSetup}
  \uses{def:physics:phyx-mini-0209:phyxminiproblems-problemphyxmini0209-lengthmagnitude, def:physics:phyx-mini-0209:phyxminiproblems-problemphyxmini0209-acousticmedium, def:physics:phyx-mini-0209:phyxminiproblems-problemphyxmini0209-troutfisherfigure}
  The acoustic setup and its unknown observables. bankSetbackX is the unknown dimensionful distance labeled x, not a preassigned answer value. Likewise, the angle fields and the predicates describing entry into the water and total reflection are independent until the governing-law premise relates them
\end{definition}
\begin{proof}
  This is the declared model definition of Trout Fisher Sound Setup.
\end{proof}
\begin{definition}[Matches Problem Data]
  \label{def:physics:phyx-mini-0209:phyxminiproblems-problemphyxmini0209-matchesproblemdata}
  \lean{PhyXMiniProblems.ProblemPhyXMini0209.MatchesProblemData}
  \uses{def:physics:phyx-mini-0209:phyxminiproblems-problemphyxmini0209-lengthinmeters, def:physics:phyx-mini-0209:phyxminiproblems-problemphyxmini0209-speedinmeterspersecond, def:physics:phyx-mini-0209:phyxminiproblems-problemphyxmini0209-troutfishersoundsetup}
  The numerical data stated in the prose: the voice is 1.8 m above the ground, and the sound speeds in air and water are respectively 343 m/s and 1440 m/s. No setback distance or answer choice occurs here
\end{definition}
\begin{proof}
  This is the declared model definition of Matches Problem Data.
\end{proof}
\begin{definition}[Matches Primary Figure]
  \label{def:physics:phyx-mini-0209:phyxminiproblems-problemphyxmini0209-matchesprimaryfigure}
  \lean{PhyXMiniProblems.ProblemPhyXMini0209.MatchesPrimaryFigure}
  \uses{def:physics:phyx-mini-0209:phyxminiproblems-problemphyxmini0209-troutfishersoundsetup}
  Primary-image evidence. The pink incident ray and its grazing refracted continuation are distinct from the dashed surface normal and the horizontal bracket x. These fields contain no numerical distance conclusion
\end{definition}
\begin{proof}
  This is the declared model definition of Matches Primary Figure.
\end{proof}
\begin{definition}[Has Physical Parameters]
  \label{def:physics:phyx-mini-0209:phyxminiproblems-problemphyxmini0209-hasphysicalparameters}
  \lean{PhyXMiniProblems.ProblemPhyXMini0209.HasPhysicalParameters}
  \uses{def:physics:phyx-mini-0209:phyxminiproblems-problemphyxmini0209-lengthinmeters, def:physics:phyx-mini-0209:phyxminiproblems-problemphyxmini0209-speedinmeterspersecond, def:physics:phyx-mini-0209:phyxminiproblems-problemphyxmini0209-troutfishersoundsetup}
  Positivity, the increase of wave speed at the boundary, and principal-angle ranges. Restricting incidence angles to [0, pi/2) selects the geometric branch shown in the diagram and makes the tangent geometry unambiguous
\end{definition}
\begin{proof}
  This is the declared model definition of Has Physical Parameters.
\end{proof}
\begin{definition}[Satisfies Critical Acoustic Refraction Physics]
  \label{def:physics:phyx-mini-0209:phyxminiproblems-problemphyxmini0209-satisfiescriticalacousticrefractionphysics}
  \lean{PhyXMiniProblems.ProblemPhyXMini0209.SatisfiesCriticalAcousticRefractionPhysics}
  \uses{def:physics:phyx-mini-0209:phyxminiproblems-problemphyxmini0209-lengthinmeters, def:physics:phyx-mini-0209:phyxminiproblems-problemphyxmini0209-speedinmeterspersecond, def:physics:phyx-mini-0209:phyxminiproblems-problemphyxmini0209-troutfishersoundsetup}
  The governing physical model. For waves, Snell's law can be written sin theta\_i / v\_air = sin theta\_r / v\_water. At the critical ray, theta\_r = pi/2. Since theta\_i is measured from the vertical normal, the right triangle in the figure has tan theta\_i = x / h. Rays below the critical incidence angle enter the water interior; rays strictly above it undergo total internal reflection. None of these relations assigns a numerical value to x
\end{definition}
\begin{proof}
  This is the declared model definition of Satisfies Critical Acoustic Refraction Physics.
\end{proof}
\begin{definition}[Keeps Trout Unfrightened]
  \label{def:physics:phyx-mini-0209:phyxminiproblems-problemphyxmini0209-keepstroutunfrightened}
  \lean{PhyXMiniProblems.ProblemPhyXMini0209.KeepsTroutUnfrightened}
  \uses{def:physics:phyx-mini-0209:phyxminiproblems-problemphyxmini0209-lengthmagnitude, def:physics:phyx-mini-0209:phyxminiproblems-problemphyxmini0209-troutfishersoundsetup}
  Sound at this setback does not propagate into the water interior
\end{definition}
\begin{proof}
  This is the declared model definition of Keeps Trout Unfrightened.
\end{proof}
\begin{definition}[Is Minimum Safe Setback]
  \label{def:physics:phyx-mini-0209:phyxminiproblems-problemphyxmini0209-isminimumsafesetback}
  \lean{PhyXMiniProblems.ProblemPhyXMini0209.IsMinimumSafeSetback}
  \uses{def:physics:phyx-mini-0209:phyxminiproblems-problemphyxmini0209-lengthmagnitude, def:physics:phyx-mini-0209:phyxminiproblems-problemphyxmini0209-lengthinmeters, def:physics:phyx-mini-0209:phyxminiproblems-problemphyxmini0209-troutfishersoundsetup, def:physics:phyx-mini-0209:phyxminiproblems-problemphyxmini0209-keepstroutunfrightened}
  The setback is the safe threshold: it itself sends no sound into the water interior, every shorter setback does, and every strictly larger setback puts the modeled ray into total internal reflection
\end{definition}
\begin{proof}
  This is the declared model definition of Is Minimum Safe Setback.
\end{proof}
\begin{definition}[Answer Choice]
  \label{def:physics:phyx-mini-0209:phyxminiproblems-problemphyxmini0209-answerchoice}
  \lean{PhyXMiniProblems.ProblemPhyXMini0209.AnswerChoice}
  Labels of the four distances printed beside the problem
\end{definition}
\begin{proof}
  This is the declared model definition of Answer Choice.
\end{proof}
\begin{definition}[meters]
  \label{def:physics:phyx-mini-0209:phyxminiproblems-problemphyxmini0209-answerchoice-meters}
  \lean{PhyXMiniProblems.ProblemPhyXMini0209.AnswerChoice.meters}
  \uses{def:physics:phyx-mini-0209:phyxminiproblems-problemphyxmini0209-answerchoice}
  Metre value printed beside each answer label
\end{definition}
\begin{proof}
  This is the declared model definition of meters.
\end{proof}
\begin{definition}[recorded Answer Choice]
  \label{def:physics:phyx-mini-0209:phyxminiproblems-problemphyxmini0209-recordedanswerchoice}
  \lean{PhyXMiniProblems.ProblemPhyXMini0209.recordedAnswerChoice}
  \uses{def:physics:phyx-mini-0209:phyxminiproblems-problemphyxmini0209-answerchoice}
  The answer label recorded in the dataset metadata
\end{definition}
\begin{proof}
  This is the declared model definition of recorded Answer Choice.
\end{proof}
\begin{definition}[Matches Displayed Centimeter]
  \label{def:physics:phyx-mini-0209:phyxminiproblems-problemphyxmini0209-matchesdisplayedcentimeter}
  \lean{PhyXMiniProblems.ProblemPhyXMini0209.MatchesDisplayedCentimeter}
  \uses{def:physics:phyx-mini-0209:phyxminiproblems-problemphyxmini0209-lengthmagnitude, def:physics:phyx-mini-0209:phyxminiproblems-problemphyxmini0209-lengthinmeters, def:physics:phyx-mini-0209:phyxminiproblems-problemphyxmini0209-answerchoice}
  Agreement to the nearest displayed centimetre. A half-centimetre tolerance is appropriate because the two sound speeds are explicitly stated only as approximate values and the answer table displays two decimal places in metres
\end{definition}
\begin{proof}
  This is the declared model definition of Matches Displayed Centimeter.
\end{proof}
% --- Archon declaration topology end ---
% --- Archon physics formalization source end ---
